\section{Conclusion}

We began with a simple observation: the aggregation step that converts per-token log-probabilities into a single confidence score is a universal but unstudied design choice in zero-cost hallucination detection. The answer is that it matters by up to 12 AUROC points, and the reason is mechanistic.

The optimal aggregation function is determined by hallucination type. Recall-failure hallucinations concentrate uncertainty at a single worst-case token; minimum log-probability is maximally discriminative. Imitative-falsehood hallucinations distribute uniform confidence throughout; mean log-probability integrates the diffuse signal. This three-step mechanism is confirmed at distributional, rank-order, and AUROC levels across two model families.

Our main contributions are: (1) the first principled aggregation selection criterion for zero-cost hallucination detection, confirmed with large effect sizes across LLaMA-2-7B and Mistral-7B-v0.1; (2) direct distributional evidence for the mechanism's first step (peakedness ratio $p = 0.002$); and (3) the discovery that raw unnormalized log-probability sum is the strongest single-pass factual-recall detector, surpassing published multi-sample Semantic Entropy at one-tenth the cost (under our evaluation protocol).

Three immediate extensions are grounded in this work: NQ completion (all code exists, GPU time only), length-stratified AUROC analysis (function implemented), and adaptive aggregation (mechanism-guided, no labels required).

The question we started with has a concrete answer: the best token log-probability aggregation function for hallucination detection depends on what kind of hallucination you expect --- and that choice is now principled, not arbitrary. The token log-probability sequence contains more structure than is typically exploited, and reading that structure correctly costs nothing.
